\chapter{The Desktop}
\label{ch:desktop}

\standfirst{The window system, in full. It is small, it is consistent, and
none of it works the way anything built after 1990 does.}

\section{What is on the screen}

A grey halftone background, and windows. Each window has a label tab in its
top-left corner and no other decoration: no title bar, no close box, no resize
grip, no scroll bar you have not asked for.

Windows are \emph{scheduled}: \ct{ScheduledControllers} holds the list, and
control passes to whichever window's rectangle contains the cursor. There is
no focus to click into. Move the pointer over a window and that window is
listening.

\section{The three buttons, again}

The colours the library calls them by are in Section~\ref{sec:buttons}; this
is what they do.

\begin{center}
\small
\begin{tabular}{@{}lp{0.62\textwidth}@{}}
\toprule
left & select\\
middle & the view's menu --- what you can do to the contents\\
right & the window's menu --- what you can do to the window\\
\bottomrule
\end{tabular}
\end{center}

\textbf{Press and hold} for both menus. Drag onto an item; release to choose
it; release outside to choose nothing.

\section{The desktop menu}
\label{sec:desktopmenu}

Press and hold the middle button anywhere the background shows through. Ten
items in three groups, and every one of them is a message sent by
\ct{ScreenController} --- the controller whose view is the grey halftone, and
which exists to provide exactly this menu.

\begin{center}
\small
\begin{tabular}{@{}lp{0.62\textwidth}@{}}
\toprule
\ctp{restore display} & redraw everything. The cure for a corrupted screen\\
\ctp{exit project} & leave the current project\\
\midrule
\ctp{project} & open a Project: a separate desktop with its own windows\\
\ctp{file list} & a file browser\\
\ctp{browser} & the System Browser\\
\ctp{workspace} & a scratch pad\\
\ctp{system transcript} & where \ctp{Transcript show:} goes\\
\ctp{system workspace} & a workspace pre-filled with expressions to run\\
\midrule
\ctp{save} & write the image, keep running\\
\ctp{quit} & save and quit, quit without saving, or continue\\
\bottomrule
\end{tabular}
\end{center}

Every one of these that opens a window then asks you to frame it: press the
left button, drag out a rectangle, release. The rest of this section is what
each one actually does.

\subsection{The screen and the project}

\begin{description}

\item[\ctp{restore display}] Sends \ct{ScheduledControllers restore}, which
repaints the desktop halftone and then asks every scheduled window to display
itself again.

You will want it more often than you expect, and that is not a fault. The
display is an ordinary \ct{Form} that any code can draw on --- including a
line you just evaluated in a workspace, which is the fastest way to learn what
\ct{Form} does and also the fastest way to paint over a Browser. Nothing is
lost when this happens; the windows are still there, still listening, and
still know how to draw themselves. This is the message that asks them to.

It is also what to reach for after the window has been resized, because
\ct{restore} re-reads \ct{Display boundingBox} rather than trusting the
rectangle a window was given when it opened.

\item[\ctp{exit project}] Leaves the project you are in and re-enters the one
that holds it (Section~\ref{sec:projects}). The top project holds
\emph{itself}, so choosing this at the top level re-enters the top project:
the screen is rebuilt from the same window list and nothing is lost.

\begin{stnote}[Why that is worth saying]
In 1983 this item ended the session. \ct{Project class>>initialize} captures
\ct{ScheduledControllers} into the top project before that variable has been
filled in, and \ct{Project>>initialProject} never assigns the transcript at
all --- so exiting the top project installed \ct{nil} as the window list and
\ct{nil} as the \ct{Transcript}, and every process that then asked either of
them for anything answered \ct{doesNotUnderstand}, several hundred times a
second, with no way back.

The bootstrap runs that initializer again once both objects exist, and fills
in the transcript the 1983 method leaves out. The menu item then does what
\ct{projectHolder _ self} always said it should.
\end{stnote}

\end{description}

\subsection{The six that open something}

\begin{description}

\item[\ctp{project}] Sends \ct{ProjectView open}. You frame a window, and that
window \emph{is} the project --- a text pane you can label with whatever the
project is for. Its middle-button menu carries one item the ordinary editor
menu does not: \ct{enter}. Choose it and the screen is replaced by that
project's own window list and its own Transcript.
Section~\ref{sec:projects} has the rest.

\item[\ctp{file list}] Sends \ct{FileList open}: a pattern pane, a list of
matching files, and a contents pane. It reads files, files code in, renames,
removes, and in this system walks into directories.
Section~\ref{sec:filelist} covers it.

\item[\ctp{browser}] Sends \ct{BrowserView openOn: SystemOrganization} --- a
System Browser on the whole system, which is where nearly all of the work in
Smalltalk gets done. Chapter~\ref{ch:browser} is about it.

\item[\ctp{workspace}] Sends \ct{StringHolderView open}: a text pane with no
model behind it and nothing to accept into. Type an expression, select it,
choose \ct{do it} or \ct{print it}. Assigning to a name you never declared
makes it a workspace variable that lives as long as the window does, which is
what makes this a scratch pad rather than a series of one-shot evaluations.

\item[\ctp{system transcript}] Sends
\ct{TextCollectorView open: Transcript label: 'System Transcript'}. The
important word is \emph{view}: there is one global \ct{Transcript}, and this
opens a window onto it. Choose the item twice and you get two windows showing
the same text.

Everything the system says out loud arrives here --- your own
\ct{Transcript show:}, and the timestamp \ct{save} writes when it takes a
snapshot. Close it and the \ct{Transcript} still exists and still collects;
you have only stopped looking.

\item[\ctp{system workspace}] A workspace that opens with text already in it,
ready to select and run.

The text is this system's own. 1983's System Workspace is a page of Xerox's
prose and example expressions, and that is precisely the part of this project
that carries no licence from anybody --- so this one says what is true of
\emph{this} system instead:

\begin{st}
3 + 4 factorial
(1 to: 10) inject: 0 into: [:a :b | a + b]
Smalltalk at: #Transcript

Display extent
Display fill: (40@40 corner: 200@120) rule: 3 mask: Form black
ScheduledControllers restore

BrowserView openOn: SystemOrganization
StringHolderView open

Processor activeProcess
[Transcript show: 'from a process'; cr] fork
\end{st}

The \ct{Display fill:} line paints a rectangle straight onto the screen, and
the \ct{ScheduledControllers restore} two lines under it is how you put the
screen back --- which is the point of having them next to each other.

\end{description}

\subsection{Saving and leaving}

\begin{description}

\item[\ctp{save}] Prompts for a name, offering \ct{snapshot}, and writes the
whole image to that name with \ct{.im} on the end. Unless the name is
\ct{snapshot}, the changes file is copied to match first, so \ct{work.im}
arrives beside \ct{work.changes}. Then the system carries on running.

Answer the prompt with nothing at all and the save is cancelled. That is the
only way out of it, and it is worth knowing, because there is no button that
says cancel.

\begin{st}
./st80 -run work.im
\end{st}

is how you come back to it. Chapter~\ref{ch:keeping} is about what is in that
file and what is not.

\item[\ctp{quit}] Opens a second menu of three:

\begin{center}
\small
\begin{tabular}{@{}lp{0.62\textwidth}@{}}
\toprule
\ctp{Save, then quit} & prompt for a name exactly as \ctp{save} does, write
  the image, and exit\\
\ctp{Quit, without saving} & exit now. Everything since the last snapshot is
  gone\\
\ctp{Continue} & never mind\\
\bottomrule
\end{tabular}
\end{center}

Releasing the button outside the menu chooses nothing, which is the same as
\ct{Continue}. An empty answer to the name prompt cancels the quit as well as
the save, so a slip at that prompt leaves you where you were rather than
exiting unsaved.

\end{description}

\section{The window menu}

The right button, anywhere over a window:

\begin{center}
\small
\begin{tabular}{@{}lp{0.62\textwidth}@{}}
\toprule
\ctp{under} & send this window behind the ones it overlaps\\
\ctp{move} & pick it up. It then follows the pointer \emph{with no button
  held}, and you click to drop it\\
\ctp{frame} & give it a new rectangle, the way opening did\\
\ctp{collapse} & shrink it to its label tab, which still has this menu\\
\ctp{close} & close it\\
\bottomrule
\end{tabular}
\end{center}

\ct{close} asks the model first, so a Browser with unaccepted edits will
interrupt and ask whether you meant it.

\section{The text editor}
\label{sec:texteditor}

Every text pane in the system is the same editor, and it has the same
middle-button menu:

\begin{center}
\small
\begin{tabular}{@{}lp{0.60\textwidth}@{}}
\toprule
\ctp{again} & repeat the last replacement, at the next occurrence\\
\ctp{undo} & undo the last change. One level\\
\midrule
\ctp{copy} \ctp{cut} \ctp{paste} & the system's clipboard, both ways\\
\midrule
\ctp{do it} & evaluate the selection, discard the answer\\
\ctp{print it} & evaluate the selection and insert the answer after it\\
\midrule
\ctp{accept} & commit --- compile the method, save the text\\
\ctp{cancel} & throw the edits away and reload\\
\bottomrule
\end{tabular}
\end{center}

In a Browser's code pane there are three more: \ct{format} (pretty-print),
\ct{spawn} (open a new browser on just this method) and \ct{explain} (say what
the selected token is).

\subsection{The clipboard is the system's}

\ct{copy} and \ct{cut} put the selection on the system's clipboard as well
as the image's own, and \ct{paste} takes the system's when it holds
something the image did not put there itself; so text moves between a
workspace and a terminal, an editor or a browser in both directions, and
the last thing copied anywhere is what pastes. Line ends are translated on
the way---the image's CR to the world's LF and back---so a method pasted
from an editor compiles, and one copied out reads as lines. Nothing else
is: the image's strings are bytes and the clipboard is UTF-8, so a
character outside ASCII pasted in arrives as the bytes of its encoding, a
curly quote as three of them. Copying inside the image and pasting inside
it keeps emphasis; what comes from outside is plain. The 1983 editor's
clipboard was a class variable that nothing outside the image could see;
\ct{lib/Clipboard} is the bridge, and there is none under \ct{-serve},
where there is no window.

\subsection{Selecting}

Press the left button and drag. Two shortcuts worth having:

\begin{itemize}
\item Double-clicking inside a pair of delimiters---quotes, brackets,
  parentheses---selects everything between them. This is how you select a
  whole block or a whole string quickly.
\item Clicking just inside the leading delimiter and dragging selects to the
  matching one.
\end{itemize}

\subsection{No keyboard shortcuts}

There is no \ct{Ctrl-D}, no \ct{Ctrl-S}, no \ct{Ctrl-Z}. \ct{do it},
\ct{accept} and \ct{undo} are menu items. The only control keys bound are a
few editing and emphasis ones: \ct{Ctrl-W} deletes the previous word,
\ct{Ctrl-B} bolds, \ct{Ctrl-I} italicises, \ct{Ctrl-X} clears emphasis.

\subsection{Moving the insertion point}

The keys a keyboard has for it all move it, and holding \ct{Ctrl} widens the
step.

\begin{center}
\small
\begin{tabular}{@{}lp{0.62\textwidth}@{}}
\toprule
\ctp{Left} \ctp{Right} & a character\\
\ctp{Up} \ctp{Down} & a line, keeping the column\\
\ctp{Ctrl-Left} \ctp{Ctrl-Right} & a word, landing on its first character
  either way\\
\ctp{Home} \ctp{End} & the ends of the line\\
\ctp{Ctrl-Home} \ctp{Ctrl-End} & the ends of the text\\
\bottomrule
\end{tabular}
\end{center}

A word is a run of letters, digits and colons, so \ct{at:put:} is one word
and not three. The line \ct{Home} and \ct{End} mean is the line you can see:
text too long for the pane wraps, and each wrapped part has its own ends.
Pressing any of these while text is selected drops the selection at the end
the key points at.

\begin{stnote}
Backspace deletes the character behind the insertion point and \ct{Delete}
the one in front of it; \ct{Return} inserts a carriage return. Typing with a
selection active replaces the selection, and either delete key removes it.
\end{stnote}

\section{Accept means compile}

In a Browser's code pane, \ct{accept} compiles the text as a method of the
selected class and installs it \emph{immediately}. There is no build step,
nothing to restart, and no way to have a stale version running.

If it does not compile you get a \ct{SyntaxError} window with the error marked
in the text. Fix it there and \ct{accept} again, or \ct{cancel}.

\begin{stwarn}
A window whose text you have edited but not accepted is holding your only copy
of that text. Closing it asks first, but \ct{cancel} does not. Accept early.
\end{stwarn}

\section{Projects}
\label{sec:projects}

\ct{project} on the desktop menu opens a \ct{Project}: a separate desktop with
its own window list, its own Transcript and its own changes. It is a virtual
desktop, from 1983, and it is useful when you want a clean screen for one task
without closing what you had.

The window you frame is a text pane, and what makes it a door rather than a
note is one extra item on its middle-button menu:

\begin{center}
\small
\begin{tabular}{@{}lp{0.62\textwidth}@{}}
\toprule
\ctp{enter} & replace the current project with this one\\
\bottomrule
\end{tabular}
\end{center}

Choose it and the screen goes empty --- not because anything closed, but
because you are looking at a different window list. \ct{exit project} on the
desktop menu comes back. Type into the pane before you enter and you have
labelled the door with what is behind it, which is the whole of the interface
the feature was given.

\section{When the screen goes wrong}

\ct{restore display} on the desktop menu redraws everything from
\ct{ScheduledControllers}. Because the display is an ordinary Form that
anything can draw on---including your own code from a workspace---this is a
normal thing to need, not a sign of a fault.

\section{Scripting the interface}

For tests and screenshots, input can be injected from the command line rather
than typed:

\begin{sh}
./st80 -inject "m 320 240; w 20; d 129; w 40; m 320 271; w 40; u 129; w 60;
                m 60 60; w 20; d 130; w 20; m 900 1400; w 20; u 130; w 200" \
       -screenshot /tmp/browser.pbm -run st80.image
\end{sh}

\ct{m X Y} moves the pointer, \ct{d} and \ct{u} press and release a button
(128 right, 129 middle, 130 left), \ct{k} types a key, \ct{w N} waits N bytecode
slices so the image has time to respond, and \ct{x} posts a window-exposed
event. That script opens the desktop menu, chooses \ct{browser}, frames a
window, and screenshots the result.
